% !TeX spellcheck = en_US
\documentclass[french]{yLectureNote}

\title{Introduction à la termodynamique}
\subtitle{Thermodynamique}
\author{Paulhenry Saux}
\date{\today}
\yLanguage{Français}

\professor{}

\usepackage{graphicx}%-pour mettre des images
\usepackage[utf8]{inputenc}%encodage
\usepackage{geometry}%pour modifier les tailles et mettre a4paper
%\usepackage{awesomebox}%pour les boites d'exercices, de pbq et de croquis d\'esactiv\'e pour les TP de PC
\usepackage{tikz}%pour deiffner + d\'ependance de chemfig
\usepackage{tkz-tab}
\usepackage{tabularx}%pour dimensionner automatiquement les tableaux avec variable X
\usepackage{awesomebox}%Pour les boites info, danger et autres
\usepackage{menukeys}%Pour deiffner les touches de Calculatrice
\usepackage{fancyhdr}%pour les en-t\^ete personnalis\'ees
\usepackage{blindtext}%pour les liens
\usepackage{hyperref}%pour les liens (\`a mettre en dernier)
\usepackage{caption}%pour la francisation de la l\'egende table vers Tableau
\usepackage{pifont}
\usepackage{array}%pour les tableaux
\usepackage{lipsum}
\usepackage{yFlatTable}
\usepackage{multicol}
\newcommand{\Lim}[1]{\lim\limits_{\substack{#1}}\:}
\renewcommand{\vec}{\overrightarrow}
\newcommand{\bz}{\overline{z}}
\DeclareMathOperator{\card}{card}
\renewcommand{\vec}{\overrightarrow}
\newcommand{\N}[0]{\mathbb{N}}
\newcommand{\R}[0]{\mathbb{R}}
\newcommand{\C}[0]{\mathbb{C}}
\newcommand{\dd}[0]{\mathrm{d}}
\newcommand{\er}{\vec{e_{\rho}}}
\newcommand{\ep}{\vec{e_{\varphi}}}
\newcommand{\pip}{\pi^+}
\newcommand{\pim}{\pi^-}
\newcommand{\mv}{\mathcal{V}}
\DeclareMathOperator{\Div}{Div}
\newcommand{\rot}{\vec{\mathrm{rot}}}
\newcommand{\grad}{\vec{\mathrm{grad}}}
\begin{document}
\setcounter{chapter}{4}

\chapter{Transition de phase - États d'équilibre}

% \checkInfo{Contexte}{Ce cours aborde l'étude des transitions de phase d'un corps pur sous l'angle de la thermodynamique classique.}

% \tipsInfo{Structure du cours}{
% \begin{itemize}
%     \item Discussion introductive (phénoménologie et expériences).
%     \item Étude de la stabilité de l'état d'équilibre.
%     \item Introduction des potentiels thermodynamiques ($G$ et $F$).
% \end{itemize}
% }



%%### Diapositive 3 : Diversité et Définitions
\section{Introduction}
\subsection{Phases et transitions}

\begin{definition}[Phase]
Ensemble d'états accessibles à un corps homogène répondant à un même jeu de lois de comportement (phénoménologie).
\end{definition}
Ce terme n'est pas strictement synonyme d'état de la matière.\marginInfo{Un système peut avoir deux phases pour un même état solide (ex: graphite et diamant sont deux phases du carbone solide).}

\begin{definition}[Transition de phase]
Changement de phénoménologie au-delà d'un seuil d'une grandeur physique fixée. 
\end{definition}

\checkInfo{Vocabulaire des changements d'état}{
\begin{itemize}
    \item \textbf{Fusion / Solidification} : Solide $\leftrightarrow$ Liquide.
    \item \textbf{Vaporisation / Liquéfaction} : Liquide $\leftrightarrow$ Gaz/Vapeur.
    \item \textbf{Sublimation / Condensation} : Solide $\leftrightarrow$ Gaz.
\end{itemize}
}



%%### Diapositive 4 : Dispositifs Expérimentaux


%\tipsInfo{Coexistence}{Dans le cas du volume imposé, la quantité de vapeur diminue continûment au profit du liquide sous l'effet de la gravité jusqu'à disparition totale de la phase gazeuse.}





%%### Diapositive 8 : Nécessité d'une approche formelle
\section{Stabilité de l'état d'équilibre}
\subsection{Approche Thermodynamique de Gibbs}

\begin{proposition}
Les conditions de transition de phase sont dictées par la stabilité de l'état d'équilibre:
\begin{itemize}
    \item \textbf{Transition possible} : Dès que l'état d'équilibre devient \textbf{métastable} (moins stable qu'un autre état accessible dans les mêmes conditions).
    \item \textbf{Transition certaine} : Si l'état d'équilibre devient \textbf{instable}.
\end{itemize}
\end{proposition}

%\tipsInfo{Lien avec l'évolution}{Cette approche permet de lier la transition de phase à la notion d'évolution spontanée du système vers un état plus stable.}


%%### Diapositive 10 : Stabilité de l'équilibre

%### Page 1 (P11) : La Métastabilité et Cas $N=1$
\subsection{Stabilité et Métastabilité}

\begin{definition}
Un état est dit \textbf{métastable} si le système est stable face à des perturbations infinitésimales, mais qu'une évolution spontanée (éloignement définitif de l'état initial) se produit au-delà d'un certain seuil de perturbation $\Delta V_0^P$.
\end{definition}


\begin{itemize}
    \item \textbf{Rompre la métastabilité} : Apporter l'énergie nécessaire pour franchir le seuil $\Delta V^P > \Delta V_0^P$\marginInfo{L'état stable est comparable à une bille dans une cuvette infinie, tandis que l'état métastable correspond à une cuvette finie (risque de basculement).}.
    \item \textbf{Supprimer la métastabilité} : Modifier les conditions pour que le seuil tende vers zéro ($(\Delta V)_0^P \to 0$), rendant l'état instable.
\end{itemize}





%### Page 2 (P12) : Stabilité en Thermodynamique (Cas $N=2$)
\subsection{Critère de stabilité thermodynamique}

\begin{definition}[État stable]
En thermodynamique, un état d'équilibre est \textbf{STABLE} si les évolutions spontanées à partir de cet état sont \textbf{impossibles}. 
\end{definition}


%### Page 3 (P13) : Synthèse des Définitions
\subsection{Classification des états d'équilibre}

\begin{definition}
Soit un état $A$ accessible pour un jeu de contraintes $(T, P \text{ ou } V)$ :
\begin{itemize}
    \item \textbf{INSTABLE} : Une perturbation infinitésimale suffit à déclencher une évolution spontanée vers un autre état.
    \item \textbf{MÉTASTABLE} : L'évolution spontanée est impossible pour de petites perturbations mais devient possible au-delà d'un seuil fini.
    \item \textbf{STABLE} : L'évolution spontanée vers un autre état accessible est rigoureusement impossible.
\end{itemize}
\end{definition}


%### Page 5 (P15) : Le Potentiel Thermodynamique
\subsection{Approche formelle par les potentiels}

\begin{definition}
Un \textbf{potentiel thermodynamique} est une fonction d'état dont les variations permettent de déterminer les possibilités d'évolution spontanée. Elle atteint sa valeur \textbf{minimale} à l'équilibre stable.
\end{definition}


%### Page 6 (P16) : Enthalpie Libre $G$ (Évolution à $T$ et $P$ fixés)
\subsection{Fonction de Gibbs}


Pour une évolution au contact d'une atmosphère ($T_0, P_0$), on définit la fonction :
$$G^* = U + P_0V - T_0S$$
La variation donne $\Delta G^* = -T_0 S_{p}$ où $S_p$ est l'entropie produite.

%     \begin{flalign*}
% G^* &= U + P_0 V - T_0 S \\
% \mathrm{d}G^* &= \mathrm{d}U + P_0\,\mathrm{d}V - T_0\,\mathrm{d}S \\
% \mathrm{d}U &= \delta Q - P_0\,\mathrm{d}V \\
% \Rightarrow\ \mathrm{d}G^* &= \delta Q - T_0\,\mathrm{d}S \\
% \mathrm{d}S &= \mathrm{d}S_e + \mathrm{d}S_p\\
% \mathrm{d}S_e &= \frac{\delta Q}{T_0} \\
% \Rightarrow \mathrm{d}S &= \frac{\delta Q}{T_0} + \mathrm{d}S_p \\
% \Rightarrow \mathrm{d}G^* &= \delta Q - T_0\left(\frac{\delta Q}{T_0} + \mathrm{d}S_p\right)\\
% &= -T_0\,\mathrm{d}S_p \\
% \Rightarrow \Delta G^* &= -T_0 S_p \le 0
% \end{flalign*}



\begin{proposition}
Au cours d'une évolution spontanée sous contraintes extérieures constantes ($T_0, P_0$), la fonction $$G^* = U + P_0V - T_0S$$ ne peut que \textbf{décroître}.
\end{proposition}

%### Page 7 (P17) : Condition de Stabilité pour $G$
\subsection{Potentiel $G$ et Équilibre stable}

\begin{theorem}
Pour des états d'équilibre à mêmes $T$ et $P$, le potentiel pertinent est l'Enthalpie Libre :
$$G = H - TS$$
L'évolution spontanée se traduit par $\Delta G < 0$.
\end{theorem}

Un état d'équilibre est \textbf{STABLE} à $T$ et $P$ fixés si $G$ y est minimale. Cela signifie qu'aucune variation (virtuelle) de $G$\marginTips{L'enthalpie libre G représente l'énergie "utile" (travail autre que celui des forces de pression) récupérable à $T$ et $P$ fixées.} ne peut être négative.



%### Page 8 (P18) : Énergie Libre $F$ (Évolution à $T$ et $V$ fixés)
\subsection{Fonction de Helmholtz}

\checkInfo{Dérivation à $V$ constant}{
Pour une évolution isochore au contact d'un thermostat ($T_0$), la variation d'énergie libre $F = U - TS$ vérifie :
$$\Delta F = -T_0 S_{p} \Rightarrow \Delta F < 0$$.
}

\begin{proposition}
L'énergie libre $F$ est le potentiel associé aux évolutions spontanées entre deux états d'équilibre où $T$ et $V$ sont identiques.
\end{proposition}


%### Page 9 (P19) : Synthèse des Critères
\subsection{Critères de Stabilité}
\begin{itemize}
    \item \textbf{À $T$ et $P$ constants} : L'état est instable/métastable si $\Delta G_{AB} < 0$ ($G_A > G_B$). Il est stable si $\Delta G_{AB} \ge 0$ ($G_A \le G_B$).
    \item \textbf{À $T$ et $V$ constants} : L'état est instable/métastable si $\Delta F_{AB} < 0$ ($F_A > F_B$). Il est stable si $\Delta F_{AB} \ge 0$ ($F_A \le F_B$).
\end{itemize}

\begin{theorem}[Critère de Gibbs-Duhem]
À partir de tout état d'équilibre stable, on a la condition générale :
$$\Delta U + P\Delta V - T\Delta S \ge 0$$
\end{theorem}


\end{document}


